\section{Related Work}
\label{sec:related}

\paragraph{Continual Graph Learning.}
Continual learning aims to sequentially learn from a stream of tasks without catastrophically forgetting prior knowledge~\citep{Kirkpatrick2017}.
In the knowledge graph domain, LKGE~\citep{Cui2023} introduced a lifelong KG embedding framework that applies EWC, experience replay, and knowledge distillation to TransE, DistMult, and RotatE~\citep{Bordes2013, Yang2015, Sun2019} on synthetic splits of FB15k-237 and WN18RR.
PS-CKGE~\citep{Wu2024} proposed pseudo-siamese networks for continual KGE, while IncDE~\citep{Song2023} addressed inductive continual KGE with decoupled embeddings.
Synaptic Intelligence~\citep{Zenke2017} offers an alternative to EWC by accumulating importance online.
However, all existing CGL methods are evaluated on synthetically partitioned generic KGs, leaving their effectiveness on real-world biomedical KG evolution untested.

\paragraph{Biomedical Knowledge Graphs.}
PrimeKG~\citep{Chandak2023} is a precision medicine KG that integrates 20+ databases into a unified schema with 129K+ nodes across 10 biomedical entity types (genes/proteins, diseases, drugs, biological processes, molecular functions, cellular components, anatomies, pathways, phenotypes, and side effects).
TxGNN~\citep{Huang2024} leverages PrimeKG for therapeutic use prediction via graph neural networks, demonstrating the value of rich biomedical KGs for drug repurposing.
Hetionet~\citep{Himmelstein2017} is an earlier integrative biomedical network that inspired PrimeKG's design.
While these resources have enabled significant advances in biomedical AI, none has been studied in a continual learning setting with real temporal evolution.

\paragraph{Temporal Knowledge Graphs.}
Temporal KGs such as ICEWS~\citep{GarciaDuran2018} and YAGO model knowledge with explicit timestamps on facts.
These are fundamentally \emph{event-based} KGs where each triple has temporal metadata (e.g., ``Obama visited France on 2015-03-01''), and the learning task is temporal link prediction or event forecasting.
In contrast, biomedical KGs undergo \emph{structural evolution}: entire entities and relation types are added or removed as scientific understanding advances, without explicit timestamps on individual facts.
\ours captures this structural evolution paradigm, which is more representative of how real-world biomedical knowledge changes over time.

\paragraph{CGL Benchmarks.}
Existing CGL benchmarks are limited in several respects (\cref{tab:comparison}).
LKGE~\citep{Cui2023} partitions FB15k-237 (a subset of Freebase) into artificial time steps by random shuffling, losing any notion of real temporal dynamics.
ICEWS~\citep{GarciaDuran2018} provides real timestamps but covers political events rather than biomedical knowledge, and its temporal structure is event-based rather than evolutionary.
PS-CKGE~\citep{Wu2024} similarly operates on generic KGs with synthetic splits.
None of these benchmarks provides multimodal node features or domain-specific evaluation tasks.
\ours addresses all of these limitations by offering real temporal evolution, multimodal features, biomedical domain specificity, and multiple evaluation tracks.

\paragraph{Knowledge Graphs and Language Models.}
Recent work has highlighted the complementary nature of knowledge graphs and large language models (LLMs)~\citep{Pan2024}.
While LLMs encode vast implicit knowledge, they suffer from hallucinations and cannot be easily updated with new factual knowledge.
KGs provide structured, verifiable knowledge that can ground LLM reasoning.
Our benchmark includes a retrieval-augmented generation (RAG) baseline~\citep{Lewis2020} that combines KG retrieval with an LLM (Qwen2.5-7B) for continual KGQA, exploring this complementary relationship in the continual learning setting.
